\section*{Erd\H{o}s Problem \#267 (Round 2)}

\subsection*{1. ROUND-2 OBJECTIVE}

I pursue \textbf{(C) obstruction/correction}.  Round~1 treated only divisibility chains such as $n_k=m^k$.  In this round I correct an omission: the case $n_{k+1}/n_k\ge 2$ is already known in the literature, and in fact it has a short self-contained proof that does \emph{not} require divisibility.  I will prove a slightly more general sufficient condition and deduce the whole $c\ge 2$ range.

\subsection*{2. Round-1 FOUNDATION USED}

I use only the most basic Round~1 ingredients:

\begin{enumerate}
\item the notation $S=\sum_k 1/F_{n_k}$;
\item the idea of comparing a rational approximation denominator against the tail of the series;
\item the standard Binet-growth heuristic for Fibonacci numbers.
\end{enumerate}

I will \emph{not} reuse the Round~1 divisibility-chain lemma; the new proof avoids divisibility entirely.

\subsection*{3. NEW INSIGHT / TOOL (ROUND-2)}

The new observation is that one can multiply the $N$th partial sum by the full product
\[
Q_N:=\prod_{j=1}^{N}F_{n_j}
\]
instead of by a nested denominator.  If the indices are sufficiently lacunary, then the tail
\[
R_N:=\sum_{k>N}\frac1{F_{n_k}}
\]
is so small that $Q_N R_N\to 0$.  This yields irrationality immediately.

This gives the following general criterion.

\medskip
\noindent
\textbf{Criterion.}
If
\[
 n_{N+1}-\sum_{j=1}^{N} n_j + N \to \infty,
\]
then $\sum_k 1/F_{n_k}$ is irrational.

The Erd\H{o}s hypothesis with $c\ge 2$ is an immediate corollary.

\subsection*{4. ATTACK PLAN (ROUND-2)}

The exact gap after Round~1 was the lack of any result beyond divisibility chains.  I now prove:

\medskip
\noindent
\textbf{Claim A.}
If $n_{N+1}-\sum_{j\le N}n_j+N\to\infty$, then $\sum_k 1/F_{n_k}$ is irrational.

\medskip
\noindent
\textbf{Claim B.}
If $n_{k+1}/n_k\ge c\ge 2$ for all $k$, then the hypothesis of Claim~A holds; hence the whole range $c\ge 2$ is settled.

This overcomes the main Round~1 obstacle because it eliminates the need for denominator divisibility.

\subsection*{5. WORK (ROUND-2)}

Let
\[
S:=\sum_{k=1}^{\infty}\frac{1}{F_{n_k}},
\qquad
S_N:=\sum_{k=1}^{N}\frac{1}{F_{n_k}},
\qquad
R_N:=S-S_N.
\]
Also define
\[
Q_N:=\prod_{j=1}^{N}F_{n_j}.
\]
Then $Q_N S_N$ is an integer.

\medskip
\noindent
\textbf{Lemma 5.1.}
For every integer $m\ge 1$,
\[
\varphi^{m-2}\le F_m\le \varphi^{m-1},
\]
where $\varphi=(1+\sqrt5)/2$.

\medskip
\noindent
\textbf{Proof.}
We argue by induction on $m$.  The bounds are true for $m=1,2$ because
\[
\varphi^{-1}\le 1\le 1,
\qquad
1\le 1\le \varphi.
\]
If they hold for $m$ and $m-1$, then using $F_{m+1}=F_m+F_{m-1}$ and $\varphi^2=\varphi+1$,
\[
F_{m+1}\le \varphi^{m-1}+\varphi^{m-2}=\varphi^m,
\]
and similarly
\[
F_{m+1}\ge \varphi^{m-2}+\varphi^{m-3}=\varphi^{m-1}.
\]
This closes the induction.
\hfill $\square$

\medskip
\noindent
\textbf{Theorem 5.2.}
Assume that
\[
\Delta_N:=n_{N+1}-\sum_{j=1}^{N}n_j+N \to \infty.
\]
Then the series
\[
\sum_{k=1}^{\infty}\frac{1}{F_{n_k}}
\]
is irrational.

\medskip
\noindent
\textbf{Proof.}
Lemma~5.1 gives
\[
Q_N=\prod_{j=1}^{N}F_{n_j}
\le \prod_{j=1}^{N}\varphi^{n_j-1}
=\varphi^{\sum_{j=1}^{N}n_j-N}.
\]
It also gives
\[
\frac{1}{F_{n_k}}\le \varphi^{2-n_k}
\qquad\mbox{for all }k.
\]
Hence
\[
R_N\le \sum_{k>N}\varphi^{2-n_k}.
\]
Since $(n_k)$ is strictly increasing, we have $n_{N+m}\ge n_{N+1}+m-1$ for $m\ge 1$, so
\[
R_N
\le \sum_{m=1}^{\infty}\varphi^{2-(n_{N+1}+m-1)}
=\varphi^{3-n_{N+1}}\sum_{m=1}^{\infty}\varphi^{-m}
\le C\varphi^{-n_{N+1}}
\]
for an absolute constant $C>0$.
Therefore
\[
Q_N R_N
\le C\varphi^{\sum_{j=1}^{N}n_j-N-n_{N+1}}
= C\varphi^{-\Delta_N}
\to 0.
\]
Now suppose, for contradiction, that $S=p/q$ is rational.  Then
\[
qQ_N R_N=qQ_N S-qQ_N S_N.
\]
The right-hand side is a positive integer, because $Q_N S_N\in\mathbf Z$ and $R_N>0$.  But the left-hand side tends to $0$, impossible.  Therefore $S$ is irrational.
\hfill $\square$

\medskip
\noindent
\textbf{Corollary 5.3.}
If
\[
\frac{n_{k+1}}{n_k}\ge 2
\qquad\mbox{for all }k,
\]
then
\[
\sum_{k=1}^{\infty}\frac{1}{F_{n_k}}
\]
is irrational.

\medskip
\noindent
\textbf{Proof.}
From $n_{k+1}\ge 2n_k$ we get, for each fixed $N$,
\[
n_j\le \frac{n_{N+1}}{2^{N+1-j}}
\qquad (1\le j\le N).
\]
Therefore
\[
\sum_{j=1}^{N}n_j
\le n_{N+1}\sum_{m=1}^{N}2^{-m}
= n_{N+1}(1-2^{-N}).
\]
Hence
\[
\Delta_N
= n_{N+1}-\sum_{j=1}^{N}n_j+N
\ge 2^{-N}n_{N+1}+N
\ge n_1+N
\to\infty.
\]
Now apply Theorem~5.2.
\hfill $\square$

\medskip
\noindent
\textbf{Corollary 5.4.}
The Erd\H{o}s question in \#267 has an affirmative answer throughout the full range $c\ge 2$.

\medskip
This recovers the known $c\ge 2$ theorem and shows exactly where the proof breaks when $1<c<2$: the product denominator $Q_N$ becomes too large compared with the first omitted Fibonacci denominator.

\subsection*{6. ADVERSARIAL VERIFICATION}

\begin{enumerate}
\item \textbf{No divisibility hidden anywhere.}  The proof uses only the product $Q_N$, never a least common multiple or a divisibility chain.
\item \textbf{Boundary case $c=2$.}  The estimate for $\Delta_N$ still tends to infinity because the extra $+N$ term remains, and in fact $\Delta_N\ge n_1+N$.
\item \textbf{Rationality contradiction.}  The expression $qQ_N R_N$ is a positive integer, not merely a rational bounded away from $0$, because $Q_N S_N$ is an integer by construction.
\item \textbf{What remains open.}  For $1<c<2$, the quantity $\Delta_N$ need not tend to infinity, so this particular product-denominator strategy genuinely stops there.
\end{enumerate}

\subsection*{7. FINAL}

\textbf{UNRESOLVED (BUT STRICTLY ADVANCED)}

\medskip
Round~2 corrects the Round~1 picture in two ways.

\begin{enumerate}
\item It identifies the already-settled range: the problem is known for $c\ge 2$.
\item It gives a short self-contained proof of the stronger criterion
\[
 n_{N+1}-\sum_{j=1}^{N}n_j+N\to\infty
 \implies
 \sum_k \frac1{F_{n_k}}\ \mbox{irrational},
\]
of which the whole range $c\ge 2$ is an immediate corollary.
\end{enumerate}

The original Erd\H{o}s problem is therefore reduced to the genuinely difficult regime
\[
1<c<2.
\]

\subsection*{8. COMPLETION ESTIMATE}

\textbf{COMPLETION: 78\%}

\subsection*{9. REFERENCES}

\begin{enumerate}
\item C. Badea, \emph{A theorem on irrationality of infinite series and applications}, Acta Arith. 63 (1993), 313--323.
\end{enumerate}
